% Named glue is calibrated using visible object bounds.
\newcommand\CoreGap[1]{\par\addvspace{#1}}
\newcommand\CoreParagraphToList{\CoreGap{5pt}}
\newcommand\CoreParagraphToTable{\CoreGap{8pt}}
\newcommand\CoreParagraphToFigure{\CoreGap{4pt}}
\newcommand\CoreParagraphToDisplayMath{\CoreGap{3pt}}
\newcommand\CoreDisplayMathToTable{\CoreGap{10pt}}
\newcommand\CoreDisplayMathToBody{\CoreGap{4pt}}
\newcommand\CoreTableToInterpretation{\CoreGap{7pt}}
\newcommand\CoreCaptionToBody{\CoreGap{5pt}}
\newcommand\CoreQuestionToNext{\CoreGap{11pt}}
\newcommand\CoreWithinGroup{\CoreGap{5pt}}
\newcommand\CoreBetweenGroups{\CoreGap{10pt}}
\AtBeginEnvironment{frame}{%
  \setlength\abovedisplayskip{5pt}\setlength\belowdisplayskip{5pt}%
  \setlength\abovedisplayshortskip{3pt}\setlength\belowdisplayshortskip{4pt}%
}
\newcommand\CoreListSpacing{%
  \setlength\itemsep{1.15pt}\setlength\parskip{0pt}%
  \setlength\parsep{0pt}\setlength\topsep{2pt}%
}
\makeatletter
\newcommand\CoreInstallListSpacing{%
  \setlength\leftmargini{19.99pt}\let\beamer@cramped\CoreListSpacing
}
\AtBeginEnvironment{itemize}{\CoreInstallListSpacing}
\AtBeginEnvironment{enumerate}{\CoreInstallListSpacing}
\AtBeginEnvironment{description}{\CoreInstallListSpacing}
\makeatother
\setbeamerfont{itemize/enumerate body}{size=\CoreBodyFont}
\setbeamerfont{itemize/enumerate subbody}{size=\CoreBodyFont}
\setbeamerfont{itemize/enumerate subsubbody}{size=\CoreBodyFont}
\setbeamertemplate{itemize item}{\textcolor{coreAccent}{\textbullet}}
\setbeamertemplate{itemize subitem}{\textcolor{coreAccent}{\textendash}}
\setbeamertemplate{itemize subsubitem}{\textcolor{coreAccent}{\textbullet}}

\newcommand\CoreSubhead[1]{\par\noindent{\usebeamerfont{core subhead}\color{coreAccent}#1}\par\vspace{2pt}}
\newcommand\TermInline[2]{\textbf{\textcolor{coreAccent}{#1}}\hspace{0.3em}#2}
\newcommand\TermBlock[2]{\CoreSubhead{#1}\noindent#2\par}
\newcommand\TeachingTakeaway[1]{\TermBlock{Takeaway}{#1}}
\newcommand\TeachingNote[1]{\TermBlock{Note}{#1}}
\newcommand\TeachingConclusion[1]{\TermBlock{Conclusion}{#1}}
\newsavebox\CoreQABox
\newcommand\CoreQALabel[1]{\textbf{\textcolor{coreAccent}{#1}}}
\newcommand\CoreQABody[1]{%
  \par\noindent%
  \sbox\CoreQABox{%
    \begin{minipage}[t]{\dimexpr\linewidth-1.1pt-.42em\relax}%
      \CoreBodyFont\raggedright\setlength\parindent{0pt}#1\par
    \end{minipage}}%
  {\color{coreAccent}\vrule width1.1pt height\ht\CoreQABox depth\dp\CoreQABox}%
  \hspace{.42em}\usebox\CoreQABox\par
}
\newcommand\TeachingQuestion[1]{%
  \CoreQABody{\CoreQALabel{Question}\hspace{.45em}{\itshape #1}}%
}
\newcommand\TeachingAnswer[1]{%
  \par\noindent\hspace*{\dimexpr1.1pt+.42em\relax}\CoreQALabel{Answer}\par
  \CoreQABody{{\upshape #1}}%
}
\newsavebox\CoreCaptionMeasure
\newsavebox\CoreCaptionLineBox
\newcount\CoreCaptionLineCount
\newcommand\CoreCountCaptionLines{%
  \let\CoreCaptionCountNext\relax
  \ifnum\lastnodetype=1\relax
    \setbox\CoreCaptionLineBox=\lastbox
    \global\advance\CoreCaptionLineCount by 1
    \let\CoreCaptionCountNext\CoreCountCaptionLines
  \else\ifnum\lastnodetype=11\relax
    \unskip\let\CoreCaptionCountNext\CoreCountCaptionLines
  \else\ifnum\lastnodetype=12\relax
    \unkern\let\CoreCaptionCountNext\CoreCountCaptionLines
  \else\ifnum\lastnodetype=13\relax
    \unpenalty\let\CoreCaptionCountNext\CoreCountCaptionLines
  \fi\fi\fi\fi
  \CoreCaptionCountNext
}
\newcommand\CoreCaption[1]{%
  \par\vspace{2pt}\begingroup
  \CoreCaptionFont\color{coreMuted}%
  \setbox\CoreCaptionMeasure=\vbox{%
    \hsize=\linewidth\raggedright\noindent#1\par}%
  \global\CoreCaptionLineCount=0
  \setbox\CoreCaptionMeasure=\vbox{%
    \unvbox\CoreCaptionMeasure\CoreCountCaptionLines}%
  \ifnum\CoreCaptionLineCount=1\relax\centering\else\raggedright\fi
  \noindent#1\par\endgroup
}
\let\ShortCaption\CoreCaption
\let\LongCaption\CoreCaption
\let\CoreTableNote\CoreCaption
\newsavebox\CoreCaptionedBox
\newcommand\CoreCaptioned[2]{%
  \sbox\CoreCaptionedBox{#1}%
  \par\noindent\begin{minipage}[t]{\wd\CoreCaptionedBox}%
    \usebox\CoreCaptionedBox\CoreCaption{#2}%
  \end{minipage}\par
}
\setbeamertemplate{caption}{\raggedright\CoreCaptionFont\color{coreMuted}\insertcaption\par}
\setlength\abovecaptionskip{2pt}
\setlength\belowcaptionskip{0pt}
\newcolumntype{Y}{>{\raggedright\arraybackslash}X}
\newcommand\CoreTableLayout{\CoreTableFont\setlength\tabcolsep{4pt}\renewcommand\arraystretch{1.19}}
\newcommand\CoreAtom[1]{\mbox{#1}}
\lstdefinestyle{presentation-code}{
  basicstyle=\CoreCodeFont, keywordstyle=\color{coreAccent}\bfseries,
  commentstyle=\color{coreMuted}, showstringspaces=false, columns=fullflexible,
  keepspaces=true, breaklines=true, breakatwhitespace=true, tabsize=2,
  frame=single, rulecolor=\color{coreRule}, backgroundcolor=\color{black!2},
  xleftmargin=.6mm, xrightmargin=.6mm, aboveskip=1pt, belowskip=.5pt,
  literate={"}{{\textquotedbl}}1 {'}{{\textquotesingle}}1
}
\lstset{style=presentation-code}

\newcommand\CoreColumns[4]{%
  \par\noindent\begin{minipage}[t]{#1\linewidth}\vspace{0pt}#3\par\end{minipage}%
  \hfill\begin{minipage}[t]{#2\linewidth}\vspace{0pt}#4\par\end{minipage}\par
}
\newcommand\CoreEqualColumns[2]{\CoreColumns{.48}{.48}{#1}{#2}}
\newcommand\CoreEvidenceLeft[2]{\CoreColumns{.57}{.39}{#1}{#2}}
\newcommand\CoreEvidenceRight[2]{\CoreColumns{.39}{.57}{#1}{#2}}

\tikzset{
  core process node/.style={draw=coreAccent,line width=.55pt,
    rounded corners=2pt,inner xsep=6pt,inner ysep=3pt,font=\CoreBodyFont,align=center},
  core process edge/.style={-{Latex[length=3pt,width=4pt]},coreAccent,line width=.65pt},
  core process distance/.style={node distance=9.5mm}
}

% Build the cue in the page background box so it does not require aux
% coordinates. Its zero-height line cannot move the ordinary body.
\newsavebox\CoreClosingCueBox
\newcommand\CoreClosingCue{%
  \sbox\CoreClosingCueBox{{\usebeamerfont{core closing question}\color{coreMuted}Questions?}}%
  \vbox to \paperheight{%
    \offinterlineskip\kern\dimexpr\headheight+.78\dimexpr\textheight+4pt\relax+.5pt\relax
    \hbox to \paperwidth{%
      \hfil\raisebox{\dimexpr(\dp\CoreClosingCueBox-\ht\CoreClosingCueBox)/2\relax}[0pt][0pt]{%
        \usebox\CoreClosingCueBox}\hfil}%
    \vfil
  }%
}
\newcommand\CoreClosing[2]{%
  \begingroup
  \setbeamertemplate{background}{\CoreClosingCue}%
  \begin{frame}[t]{#1}%
    \begin{minipage}[t][.68\textheight][t]{\textwidth}#2\par\end{minipage}%
  \end{frame}%
  \endgroup
}
\newcommand\backmatter[1][]{%
  \ifstrequal{#1}{notitle}{\CoreClosing{Discussion}{\inserttitle}}%
    {\CoreClosing{\inserttitle}{\insertsubtitle\par\insertshortinstitute}}%
}
\newenvironment{colorblock}[3][white]{%
  \begingroup\setbeamercolor{block title}{fg=#1,bg=#2}%
  \setbeamercolor{block body}{fg=#1,bg=#2}\begin{block}{#3}%
}{\end{block}\endgroup}
